\documentclass[12pt]{article}
\usepackage[utf8]{inputenc}
\usepackage[margin=0.75in]{geometry}
\usepackage{amsmath}\usepackage{amssymb}\usepackage{enumitem}\usepackage{multicol}\usepackage{tcolorbox}
\setlength{\parindent}{0pt}
\begin{document}

\begin{center}{\Large \textbf{Factor out the greatest common factor (GCF)}}\\[2pt]{\small Algebra I \textbar\ CK-12: Factoring when the Leading Coefficient Equals 1}\end{center}\vspace{0.25cm}

{\large\textbf{Key Idea}}\par\smallskip
The \textbf{GCF} is the largest factor every term shares — both the number and the variables. Factoring it out is the reverse of distributing.\par


{\large\textbf{Key Terms}}\par\smallskip
\begin{itemize}[leftmargin=1.4em]
  \item \textbf{GCF} — greatest common factor
  \item \textbf{Factor} — rewrite as a product
\end{itemize}


{\large\textbf{Worked steps}}\par\smallskip
Step 1: Find the GCF of the numbers and the smallest power of each variable.\par
Step 2: Divide each term by the GCF and write GCF(quotient).\par
Example: $6x^2+9x$: GCF is $3x$, so $6x^2+9x=3x(2x+3)$.\par


{\large\textbf{You Try}}\par\smallskip
Factor $10x+15$.\par

\end{document}